% answers_short.tex —— 2009 年·普通高等学校招生考试·湖北卷·理科 · 简短答案速查
% 由 tools/make_short_answers.py 自 answers.tex 的 \gksolhead/\gkanswer 提取，
% 并与 2009/2009_数学真题_湖北理_答案要点.md 逐题比对后生成，勿手改（改源 MD 后重跑）。

\begin{gknotice}
\textbf{使用说明}：本节只给最终答案便于快速核对，完整推导见解析版。答案为编写者独立推导，\textbf{非官方标准答案}。
\end{gknotice}

\gkpart{选择题}
\begin{center}\begin{tabular}{ccccccccccc}
\toprule
题号 & 1 & 2 & 3 & 4 & 5 & 6 & 7 & 8 & 9 & 10 \\
\midrule
答案 & A & D & C & B & C & B & A & B & D & C \\
\bottomrule
\end{tabular}\end{center}

\gkpart{填空题}
\begin{itemize}[leftmargin=2.6em]
  \item[\textbf{11.}] $-2$。
  \item[\textbf{12.}] $64$；$0.4$。
  \item[\textbf{13.}] $12800\arccos\dfrac{8}{53}$。
  \item[\textbf{14.}] $1$。
  \item[\textbf{15.}] $4$，$5$，$32$。
\end{itemize}

\gkpart{解答题}
\begin{itemize}[leftmargin=2.6em]
  \item[\textbf{16.}] $P(\eta=5)=\dfrac1{16}$，$P(\eta=6)=\dfrac18$，$P(\eta=7)=\dfrac3{16}$，$P(\eta=8)=\dfrac14$，$P(\eta=9)=\dfrac3{16}$，$P(\eta=10)=\dfrac18$，$P(\eta=11)=\dfrac1{16}$；$E\eta=8$。
  \item[\textbf{17.}] （1）$2$；（2）$\cos\beta=1$ 或 $\cos\beta=0$。
  \item[\textbf{18.}] （1）证明见解析；（2）$\sqrt2$。
  \item[\textbf{19.}] （1）$a_n=\dfrac{n}{2^n}$；（2）$n=1$，$2$ 时 $T_n<\dfrac{5n}{2n+1}$，$n\geqslant3$ 时 $T_n>\dfrac{5n}{2n+1}$。
  \item[\textbf{20.}] （1）证明见解析；（2）存在，$\lambda=4$。
  \item[\textbf{21.}] （1）$b=-1$，$c=3$；（2）$b\neq0$ 时两条切线与曲线的公共点分别为 $(0,bc)$，$(3b,4bc)$ 和 $\left(2b,\dfrac{4b^3}{3}+3bc\right)$，$\left(-b,\dfrac{4b^3}{3}\right)$；$b=0$ 时两切线重合为 $y=cx$，公共点只有 $(0,0)$；（3）$\dfrac12$。
\end{itemize}
